\documentclass[10pt,a4paper]{article}
\usepackage{amsmath,amssymb,amsthm,graphicx}
\usepackage{amsmath,amssymb,amsthm}
\newcommand{\FF}{\mathbb F}
\newcommand{\PP}{\mathbb P}
\newcommand{\Aspace}{\mathbb A}
\newcommand{\GG}{\mathrm G}
\newcommand{\supp}{\operatorname{supp}}
\newenvironment{result}[1]{\par\medskip\noindent\textbf{#1.}\itshape}{\par\medskip}
\title{Explicit variety-evasive subspace families for bounded-degree projective and affine varieties}
\author{Zeyu Guo}
\date{doi:10.1007/s00037-024-00256-1; CC BY 4.0}
\begin{document}
\maketitle
\noindent\textit{Editorial note.} Faithful editable transcription of complete actual CC BY2024 journal main construction and necessary author proofs, using the licensed61-pagePDF directly. Full characteristic-dependent coefficients/bit complexity and improved journal N(k,d,n) retained. The author parallel dual-Stiefel case wording and original printed indices are preserved without expanding or repairing mathematics. Complete directNoether applications retained in the same record; PIT application and further strengthened affine sketch are not claimed as extra tasks. Standard rank-condensing/geometry facts retain author citations without recursive reproving. Journal bibliography metadata transcribed from publisher citation\_reference fields. Task-specific files do not mutate handed-offCC065 materials. All production/extraction/validation solely on Yupeng.\par\medskip
\noindent\textit{Allowed standard prerequisites (editorial).}\begin{itemize}
\item Projective algebraic varieties, Grassmannians/Plucker coordinates, Chow embeddings and the standard dimension formula for Chow varieties. Original definitions and citedTheorems3.5/3.7/4.1/4.4 retained, with author-provided explanatory argument. Classical results and projective hyperplane dimension fact cited rather than recursively reproved.
\end{itemize}\medskip
\section*{Original ambient variety and evasion definitions (journal pp.3--6)}
Let $\FF$ be an algebraically closed field. An affine $n$-space $\Aspace^n$, as a set, is simply defined to be the vector space $\FF^n$. We also need the notion of a projective $n$-space, denoted by $\PP^n$, which is (intuitively) the set of lines passing through the origin $0$ of $\Aspace^{n+1}$. Formally, it is defined to be the quotient set $(\Aspace^{n+1}\setminus\{0\})/\sim$, where $\sim$ is the equivalence relation defined by scaling, i.e., $u\sim v$ if $u=cv$ for some nonzero scalar $c\in\FF$.
An (affine) subvariety $\mathcal V\subseteq\Aspace^n$ is the set of common zeros of a set of $n$-variate polynomials over $\FF$. Similarly, a (projective) subvariety $\mathcal V\subseteq\PP^n$ is the set of common zeros of a set of homogeneous $(n+1)$-variate polynomials over $\FF$, where we represent each element of $\PP^n$ as an $(n+1)$-tuple in $\Aspace^{n+1}$. In this paper, a variety refers to a subvariety of a projective or affine space, and is said to be irreducible if it cannot be written as a union of finitely many proper subvarieties. (Varieties in this paper are not necessarily irreducible. They are often called algebraic sets in literature.)
The dimension of a variety $\mathcal V$, denoted by $\dim(\mathcal V)$, is intuitively the ``degree of freedom'' of picking a point in the variety. See Section 2.3 for its formal definition. In particular, for a linear subspace $\mathcal V\subseteq\Aspace^n$, the dimension of $\mathcal V$ as a variety is simply its linear-algebraic dimension.
\begin{result}{Definition 1.1 (Evading)}
Let $\mathcal V_1$ and $\mathcal V_2$ be irreducible subvarieties of $\PP^n$ or $\Aspace^n$. We say $\mathcal V_1$ evades $\mathcal V_2$ if
\[\dim(\mathcal V_1\cap\mathcal V_2)\le\dim(\mathcal V_1)+\dim(\mathcal V_2)-n,\]
where the dimension of an empty set is assumed to be $-\infty$. In particular, if $\dim(\mathcal V_1)+\dim(\mathcal V_2)<n$, then $\mathcal V_1$ evades $\mathcal V_2$ iff $\mathcal V_1\cap\mathcal V_2=\emptyset$.
More generally, suppose $\mathcal V_1$ is irreducible but $\mathcal V_2$ is possibly reducible. We say $\mathcal V_1$ evades $\mathcal V_2$ if it evades every irreducible component of $\mathcal V_2$.
\end{result}
\begin{result}{Definition 1.2 (Subspace family)}
For $0\le k\le n$, a finite collection (i.e., multiset) of $k$-subspaces of $\PP^n$ is called a (projective) $k$-subspace family on $\PP^n$. Similarly, a finite collection of affine $k$-subspaces of $\Aspace^n$ is called an affine $k$-subspace family on $\Aspace^n$.
\end{result}
\begin{result}{Definition 1.3 (Variety evasive subspace family)}
Let $\mathcal F$ be a family of subvarieties of $\PP^n$ (resp. $\Aspace^n$). Let $\mathcal H$ be a $k$-subspace family on $\PP^n$ (resp. affine $k$-subspace family on $\Aspace^n$) where $0\le k\le n$. Then:
\begin{itemize}\item We say $\mathcal H$ is $\mathcal F$-evasive if for every $\mathcal V\in\mathcal F$, there exists $W\in\mathcal H$ that evades $\mathcal V$.
\item We say $\mathcal H$ is $(\mathcal F,\epsilon)$-evasive if for every $\mathcal V\in\mathcal F$, a random element $W\in\mathcal H$ evades $\mathcal V$ with probability at least $1-\epsilon$.\end{itemize}
\end{result}
\begin{result}{Definition 1.4 (degree)}
The degree of an irreducible variety $\mathcal V$ in $\PP^n$ (resp. $\Aspace^n$) is the number of intersections of $\mathcal V$ with a general projective (resp. affine) subspace of codimension $\dim(\mathcal V)$. Following Heintz (1983); Heintz \& Schnorr (1980), we define the degree of a (possibly reducible) variety to be the sum of the degrees of its irreducible components.
\end{result}
\section*{Original Pl\"ucker and Chow-form background (journal pp.27--30)}
\paragraph{Grassmannians.}
Let $n\in\mathbb N$ and $k\in\{0,1,\ldots,n-1\}$. The Grassmannian $\GG(k+1,n+1)$ is the set of all $(k+1)$-dimensional linear subspaces of $\Aspace^{n+1}$. By taking the quotient modulo scalars, it may also be identified with the set of all $k$-subspaces of $\PP^n$, which we denote by $\GG(k,n)$.
\paragraph{The Pl\"ucker embedding and Pl\"ucker coordinates.}
Consider a linear subspace $W\in\GG(k+1,n+1)$. The simplest way of representing $W$ is using a $(k+1)\times(n+1)$ matrix $A$ over $\FF$ such that $W$ equals the row space of $A$. We call such a matrix $A$ a generating matrix of $W$. For convenience, we also say $A$ is a generating matrix of $\PP W\in\GG(k,n)$.
\begin{result}{Definition 3.4 (Pl\"ucker embedding)}
Define $\phi:\GG(k+1,n+1)\to\PP^{\binom{n+1}{k+1}-1}$ by
\[\phi(W)=(\det A_{[k+1],S})_{S\in\binom{[n+1]}{k+1}}\]
where $A$ is a generating matrix of $W$.
\end{result}
The Pl\"ucker embedding embeds the Grassmannian $\GG(k+1,n+1)$ in $\PP^{\binom{n+1}{k+1}-1}$ as an irreducible projective subvariety, as stated by the following theorem. See, e.g., Fulton (1997); Harris (1992) for proofs.
\begin{result}{Theorem 3.5}
The Pl\"ucker embedding $\phi$ is a well-defined injective map whose image is an irreducible projective subvariety of $\PP^{\binom{n+1}{k+1}-1}$.
\end{result}
The homogeneous coordinates $(\det A_{[k+1],S})_{S\in\binom{[n+1]}{k+1}}$ of $\phi(W)$ are called the (primal) Pl\"ucker coordinates of $W$.
Denote by $R:=\FF[X_S:S\in\binom{[n+1]}{k+1}]$ the homogeneous coordinate ring of $\PP^{\binom{n+1}{k+1}-1}$. The irreducible projective subvariety $\phi(\GG(k+1,n+1))$ is defined by a homogeneous prime ideal of $R$, which we denoted by $I$. Then $R/I$ is the homogeneous coordinate ring of $\phi(\GG(k+1,n+1))$. The ideal $I$ contains precisely the polynomial relations that the Pl\"ucker coordinates need to satisfy. It is also known that $I$ is generated by certain quadratic forms, known as the Pl\"ucker relations. See Fulton (1997); Harris (1992) for details.
\paragraph{Dual Pl\"ucker coordinates.}
Alternatively, we could represent a linear subspace $W\in\GG(k+1,n+1)$ by an $(n-k)\times(n+1)$ matrix $B$ over $\FF$ whose rows specify the linear equations defining $W$. We call such a matrix $B$ a parity check matrix of $W$. For convenience, we also say $B$ is a parity check matrix of $\PP W\in\GG(k,n)$.
The entries of $B$ are called the dual Stiefel coordinates of $W$. This gives another embedding $\phi^\vee:\GG(k+1,n+1)\to\PP^{\binom{n+1}{n-k}-1}=\PP^{\binom{n+1}{k+1}-1}$, defined by
\[\phi^\vee(W)=(\det B_{[n-k],S})_{S\in\binom{[n+1]}{n-k}}.\]
The homogeneous coordinates $(\det B_{[n-k],S})_{S\in\binom{[n+1]}{n-k}}$ of $\phi^\vee(W)$ are called the dual Pl\"ucker coordinates of $W$.
Recall that we denote by $\GG(k,n)$ the set of all $k$-subspaces of $\PP^n$. By identifying $\GG(k+1,n+1)$ with $\GG(k,n)$ via $W\mapsto\PP W$, we regard $\phi$ and $\phi^\vee$ as maps from $\GG(k,n)$ to $\PP^{\binom{n+1}{k+1}-1}$.
\begin{result}{Definition 3.6 (Gelfand et al. 1994)}
For an irreducible subvariety $\mathcal V\subseteq\PP^n$ of dimension $n-k-1$, define the associated hypersurfaces $Z_{\mathcal V}$ of $\mathcal V$ as the set of $k$-subspaces intersecting $\mathcal V$, i.e.,
\[Z_{\mathcal V}:=\{W\in\GG(k,n):\mathcal V\cap W\ne\emptyset\}.\]
\end{result}
The term ``associated hypersurface'' is justified by the following theorem.
\begin{result}{Theorem 3.7}
Let $\mathcal V\subseteq\PP^n$ be an irreducible projective subvariety of dimension $n-k-1$ and degree $d\in\mathbb N^+$. Then there exists a nonzero homogeneous polynomial $P_{\mathcal V}\in R=\FF[X_S:S\in\binom{[n+1]}{k+1}]$ of degree $d$ such that $\phi(Z_{\mathcal V})$ is defined by $P_{\mathcal V}$ as a subvariety of $\phi(\GG(k,n))$. That is,
\[\phi(Z_{\mathcal V})=\phi(\GG(k,n))\cap V(P_{\mathcal V}).\]
Moreover, $R_{\mathcal V}:=P_{\mathcal V}+I\in(R/I)_d$ is uniquely determined by $\mathcal V$ up to scalars.
\end{result}
Theorem 3.7 is explicitly stated as Dalbec \& Sturmfels (1995, Theorem 1.1 and Corollary 2.1). A proof can be found in Gelfand et al. (1994, Section 3.2). We briefly explain how to find a polynomial $P_{\mathcal V}$ satisfying Theorem 3.7: Firstly, it can be shown using the trick of dimension counting via incidence varieties that $\phi(Z_{\mathcal V})$ is an irreducible projective subvariety of the Grassmannian $\phi(\GG(k,n))$ of codimension one (Gelfand et al. 1994, Section 3.2, Proposition 2.2). Secondly, the homogeneous coordinate ring $R/I$ of the Grassmannian is known to be a unique factorization domain (Fulton 1997, Chapter 9). These two facts imply that the homogeneous ideal of $R/I$ defining $\phi(Z_{\mathcal V})$ is a principal ideal. Choose $R_{\mathcal V}$ to be a generator of this principal ideal, which is unique up to scalars. Then lift $R_{\mathcal V}\in R/I$ to $P_{\mathcal V}\in R$.
\begin{result}{Definition 3.8 (Chow form)}
Let $\mathcal V\subseteq\PP^n$ be an irreducible subvariety of dimension $n-k-1$ and degree $d\in\mathbb N^+$. Define the Chow form of $\mathcal V$ in Pl\"ucker coordinates, or simply the Chow form of $\mathcal V$, to be $R_{\mathcal V}\in(R/I)_d$ as in Theorem 3.7.
More generally, for an equidimensional subvariety $\mathcal V=\bigcup_{i=1}^s\mathcal V_i\subseteq\PP^n$ of dimension $n-k-1$ and degree $d$, where $\mathcal V_1,\ldots,\mathcal V_s$ are the irreducible components of $\mathcal V$, the Chow form of $\mathcal V$ is $R_{\mathcal V}:=\prod_{i=1}^sR_{\mathcal V_i}\in(R/I)_d$. It is uniquely determined by $\mathcal V$ up to scalars.
\end{result}
As a $k$-subspace intersects $\mathcal V=\bigcup_{i=1}^s\mathcal V_i$ iff it intersects some $\mathcal V_i$, we see from Theorem 3.7 that the Chow form $R_{\mathcal V}$ of an equidimensional projective subvariety $\mathcal V$ of dimension $n-k-1$ vanishes precisely at the set of $k$-subspaces that intersect $\mathcal V$.

\section*{Original main statement (journal p.7)}
For $n,d\in\mathbb N^+$ and $k\in\{0,1,\ldots,n-1\}$, define $N(k,d,n)$ by
\[N(k,d,n):=\min\left\{\binom{(k'+1)(n'+1+d)}{(k'+1)d},\binom{(n-k)(n'+1+d)}{(n-k)d}\right\},\]
where $k':=\min\{k,d-2\}\le k$ and $n':=k'+n-k\le n$.
\begin{result}{Theorem 1.7 (Main Theorem)}
For $n,d\in\mathbb N^+$, $k\in\{0,1,\ldots,n-1\}$, and $\epsilon\in(0,1)$, there exists an $(n,d,\epsilon)$-evasive $k$-subspace family (resp. affine $k$-subspace family) $\mathcal H$ on $\PP^n$ (resp. $\Aspace^n$) of size
\[\operatorname{poly}(N(k,d,n),n,1/\epsilon)\le\operatorname{poly}(n^{\min\{k+1,n-k,d\}d},d,1/\epsilon).\]
Moreover:
\begin{itemize}
\item If $\operatorname{char}(\FF)=0$, where $\FF$ denotes the base field, then the linear equations defining the projective or affine subspaces in $\mathcal H$ are defined over $\mathbb Q$. Moreover, the bit-lengths of the numerators and denominators of the coefficients of these linear equations are polynomial in $|\mathcal H|$, and the total bit complexity of computing these linear equations is polynomial in $|\mathcal H|$.
\item If $\operatorname{char}(\FF)=p>0$, then the linear equations defining the projective or affine subspaces in $\mathcal H$ are defined over an extension field $\FF_q$ of $\FF_p$, where $q\le\operatorname{poly}(|\mathcal H|,p)$. The total bit complexity of computing these linear equations and constructing the field $\FF_q$ is polynomial in $|\mathcal H|$ and $\log p$.
\end{itemize}
In particular, when $d$ is bounded, the bit complexity of constructing $\mathcal H$ is polynomial in $n/\epsilon$ (and $\log p$ if $\operatorname{char}(\FF)=p>0$).
\end{result}

\section*{Original degree-bounded evasion convention (journal p.6)}
\begin{result}{Definition 1.5}
We say a projective (resp. affine) $k$-subspace family $\mathcal H$ on $\PP^n$ (resp. $\Aspace^n$) is $(n,d)$-evasive if it is $\mathcal F$-evasive, where $\mathcal F$ is chosen to be the family of all subvarieties of $\PP^n$ (resp. $\Aspace^n$) of degree at most $d$. Similarly, we say $\mathcal H$ is $(n,d,\epsilon)$-evasive if it is $(\mathcal F,\epsilon)$-evasive.
\end{result}
\begin{result}{Remark 1.6}
In Definition 1.5, we do not make any assumption about the dimension of the varieties in $\mathcal F$ or their irreducible components. We will see in Section 3.1 that in fact, it suffices to consider the subfamily of equidimensional varieties or even irreducible varieties of dimension $n-k-1$ when constructing variety evasive $k$-subspace families.
\end{result}
\section*{Original Section 2.1: Black-Box PIT for Low Degree Polynomials (journal pp.18--20)}
\begin{result}{Definition 2.1 ($\epsilon$-hitting set)}
Let $\mathcal F$ be a family of polynomials in $\FF[X_1,\ldots,X_n]$ and $\epsilon\in(0,1)$. We say a finite collections of points $\mathcal H\subseteq\FF^n$ is an $\epsilon$-hitting set for $\mathcal F$ if for any nonzero $Q\in\mathcal F$, the evaluation $Q(\alpha)$ is nonzero for all but at most $\epsilon$-fraction of $\alpha\in\mathcal H$.
\end{result}
We need an explicit construction of $\epsilon$-hitting sets for low degree polynomials. This problem has been well studied (Bl\"aser \& Pandey 2020; Bogdanov 2005; Bshouty 2014; Cohen \& Ta-Shma 2013; Demillo \& Lipton 1978; Klivans \& Spielman 2001; Lu 2012; Schwartz 1980; Zippel 1979). For completeness, we present a construction based on sparse polynomial identity testing.
Recall that a polynomial is $s$-sparse if it has at most $s$ monomials. We need the following lemma from Agrawal et al. (2015).
\begin{result}{Lemma 2.2 (Agrawal et al. 2015, Lemma 4, restated)}
For $n,s,d\in\mathbb N^+$ and $\epsilon_0\in(0,1)$, there exist maps $w_1,w_2,\ldots,w_N:[n]\to[N\log N]$, where $N=\operatorname{poly}(n,s,\log d,\epsilon_0^{-1})$, such that for any nonzero $s$-sparse polynomial $f\in\FF[X_1,\ldots,X_n]$ of individual degree at most $d$, all but at most $\epsilon_0$-fraction of $w_i$ among $w_1,\ldots,w_N$ satisfies $f(Y^{w_i(1)},\ldots,Y^{w_i(n)})\ne0$. Moreover, the bit complexity of computing $w_1,w_2,\ldots,w_N$ is polynomial in $N$.
\end{result}
Given $n,d\in\mathbb N^+$ and $\epsilon\in(0,1)$, we construct an $\epsilon$-hitting set for $n$-variate polynomials of degree at most $d$ as follows:
\begin{enumerate}
\item Let $s=\binom{n+d}{d}$, $\epsilon_0=\epsilon/2$, and $M=\lceil\epsilon_0^{-1}dN\log N\rceil$, where $N$ is as in Lemma 2.2.
\item Let $w_1,\ldots,w_N$ be as in Lemma 2.2, which can be computed in time $\operatorname{poly}(N)$.
\item If $\operatorname{char}(\FF)=0$, let $S=[M]\subseteq\mathbb Z\subseteq\FF$. If $\operatorname{char}(\FF)=p>0$, choose the smallest $p$-power $q$ such that $q\ge M$, and choose $S$ to be a subset of $\FF_q\subseteq\FF$ of cardinality $M$. We remark that $\FF_q$ can be constructed deterministically in time $\operatorname{poly}(M,\log p)$. To see this, note that $q\le Mp$ by the minimality of $q$. If $M\le p$, then $\FF_q$ is just $\FF_p$. On the other hand, if $p<M\le q$, then $\FF_q$ can be constructed in time $\operatorname{poly}(p,[\FF_q:\FF_p])$ (see, e.g., Lenstra (1990)), which is polynomial in $M$ since $p<M$ and $q\le Mp$.
\item Finally, construct the following collection of points in $\FF^n$ of size $MN$
\[T=\{(\alpha^{w_i(1)},\ldots,\alpha^{w_i(n)}):\alpha\in S,i\in[N]\}\subseteq\FF^n.\]
\end{enumerate}
\begin{result}{Lemma 2.3}
For any nonzero polynomial $f\in\FF[X_1,\ldots,X_n]$ of degree at most $d$, we have $f(u)\ne0$ for all but at most $\epsilon$-fraction of $u\in T$. The collection $T$ has cardinality $\operatorname{poly}(\binom{n+d}{d},1/\epsilon)$ and can be computed in time $\operatorname{poly}(|T|)$.
\end{result}
\begin{proof}
Let $f\in\FF[X_1,\ldots,X_n]$ be a nonzero polynomial of degree at most $d$. Note that $f$ is trivially $s$-sparse, where $s=\binom{n+d}{d}$. So by Lemma 2.2, for all but at most $\epsilon_0$-fraction of $i\in[N]$, we have $\widetilde f_i:=f(Y^{w_i(1)},\ldots,Y^{w_i(n)})\ne0$. Consider $i\in[N]$ such that $\widetilde f_i\ne0$. Note that $\widetilde f_i$ is a univariate polynomial of degree at most $dN\log N$. So it has at most $dN\log N\le\epsilon_0M$ zeros. Therefore, by the choice of $M$, we have $f(\alpha^{w_i(1)},\ldots,\alpha^{w_i(n)})=\widetilde f_i(\alpha)\ne0$ for all but at most $\epsilon_0$-fraction of $\alpha\in S$. It follows that $f(u)\ne0$ holds for all but at most $\epsilon$-fraction of $u\in T$, as claimed. The rest of the lemma follows easily from the construction.
\end{proof}
Note that the seed length required to choose a random element in $T$ is $\log|T|=O(\log\binom{n+d}{d}+\log(1/\epsilon))$, which is optimal up to a constant factor. We have made no effort to optimize the constant hidden in $O(\cdot)$. Interested readers may find the state-of-the-art result in Bl\"aser \& Pandey (2020), which achieves the optimal constant, at least when $d=n^{o(1)}$.
\section*{Original Section 2.2: Explicit Lossless Rank Condensers (journal pp.20--21)}
We need the following lemma in the context of lossless rank condensers. The construction in the lemma was given by Forbes \& Shpilka (2012) and the lemma itself follows implicitly from the analysis of Forbes et al. (2014). It was also stated explicitly as Forbes (2014, Theorem 5.4.3).
\begin{result}{Lemma 2.4 (Forbes 2014; Forbes et al. 2014)}
Let $n\in\mathbb N^+$ and $r\in[n]$. Let $\omega\in\FF^\times$ such that the multiplicative order of $\omega$ is at least $n$. Define the $r\times n$ matrix $W=(w_{i,j})_{i\in[r],j\in[n]}$ over $\FF[X]$ by
\[w_{i,j}=(\omega^{i-1}X)^{j-1}.\]
Then for every $n\times r$ matrix $M$ over $\FF$ of rank $r$, the polynomial $\det(WM)\in\FF[X]$ is nonzero and has degree at most $r(n-r)$ after dividing out powers of $X$.
\end{result}
\begin{result}{Corollary 2.5}
Let $n,r,W$ be as in Lemma 2.4 and $\epsilon\in(0,1)$. Let $S\subseteq\FF^\times$ be a finite set of cardinality at least $r(n-r)/\epsilon$. For every $n\times r$ matrix $M$ over $\FF$ of rank $r$, we have $\operatorname{rank}(W(\alpha)M)=r$ for all but at most $\epsilon$-fraction of $\alpha\in S$, where $W(\alpha)$ denotes the matrix $(w_{i,j}(\alpha))_{i\in[r],j\in[n]}$ over $\FF$.
\end{result}
Corollary 2.5 states that the collection $\{W(\alpha):\alpha\in S\}$ of matrices is a (weak) $(r,\epsilon|S|)$-lossless rank condenser, as defined in Forbes \& Guruswami (2015). Note that for each $\alpha\in S$, we have $\operatorname{rank}(W(\alpha))=r$ and hence $W(\alpha)$ corresponds to an $(r-1)$-subspace $U_{W(\alpha)}$ of $\PP^{n-1}$. As explained in Introduction, the collection $\mathcal H=\{U_{W(\alpha)}:\alpha\in S\}$ is an $(\mathcal F,\epsilon)$-evasive $(r-1)$-subspace family on $\PP^{n-1}$, where $\mathcal F$ is the family of $(n-r-1)$-subspaces of $\PP^{n-1}$. Choosing $S$ of size $r(n-r)+1$ and $\epsilon=1-\frac1{r(n-r)+1}$ shows that the lower bound in Theorem 1.10 is achieved when $d=1$.

\section*{Original geometry background needed for the construction (journal pp.23--25)}
\paragraph{Morphisms.}
Let $\mathcal V_1\subseteq\Aspace^n$ and $\mathcal V_2\subseteq\Aspace^m$ be affine varieties. A morphism from $\mathcal V_1$ to $\mathcal V_2$ is a map $f:\mathcal V_1\to\mathcal V_2$ that is a restriction of a polynomial map $\Aspace^n\to\Aspace^m$. Such a morphism $f$ is associated with a ring homomorphism $f^\sharp:\FF[\mathcal V_2]\to\FF[\mathcal V_1]$, making $\FF[\mathcal V_1]$ an algebra over $\FF[\mathcal V_2]$. We say $f$ is finite if $\FF[\mathcal V_1]$ is finitely generated as an $\FF[\mathcal V_2]$-module.
Let $f:\mathcal V_1\to\mathcal V_2$ be a map between projective varieties $\mathcal V_1$ and $\mathcal V_2$. We say $f$ is a morphism from $\mathcal V_1$ to $\mathcal V_2$ if there exists a collection of open subsets $\{U_i\}_{i\in I}$ of $\mathcal V_2$ such that $\mathcal V_2=\bigcup_{i\in I}U_i$ (i.e., $\{U_i\}_{i\in I}$ is an open cover of $\mathcal V_2$) and for each $i\in I$, the restriction $f|_{f^{-1}(U_i)}:f^{-1}(U_i)\to U_i$ is a morphism between affine varieties. Furthermore, if each $f|_{f^{-1}(U_i)}$ is finite, then we say $f$ is finite. Finiteness does not depend on the choice of the affine open cover. Namely, if $f:\mathcal V_1\to\mathcal V_2$ is a finite morphism between projective varieties $\mathcal V_1$ and $\mathcal V_2$, and $U$ is an open subset of $\mathcal V_2$ such that $f|_{f^{-1}(U)}:f^{-1}(U)\to U$ is a morphism between affine varieties, then $f|_{f^{-1}(U)}$ is also finite.
The image of a morphism $f:\mathcal V_1\to\mathcal V_2$ is denoted by $\operatorname{Im}(f)$ or $f(\mathcal V_1)$. The image of a closed set under a finite morphism is still closed. The composition of two finite morphisms is still finite.
\paragraph{Dimension.}
The dimension of an irreducible variety $\mathcal V$, denoted by $\dim(\mathcal V)$, is the largest integer $m$ such that there exists a chain of irreducible varieties $\emptyset\subsetneq\mathcal V_0\subsetneq\mathcal V_1\subsetneq\cdots\subsetneq\mathcal V_m=\mathcal V$. More generally, the dimension of a nonempty variety is the maximal dimension of its irreducible components. We define the dimension of an empty set to be $-\infty$. A variety is equidimensional if its irreducible components have the same dimension.
If $\pi:\mathcal V\to\mathcal V'$ is a finite morphism, then $\dim(\mathcal V)=\dim(\pi(\mathcal V))$.
\paragraph{Projective closure.}
The affine $n$-space $\Aspace^n$ may be regarded as an open subset of $\PP^n$ via the map $(x_1,\ldots,x_n)\mapsto(x_1,\ldots,x_n,1)$. The complement $H_\infty:=\PP^n\setminus\Aspace^n$ is a hyperplane of $\PP^n$ defined by $X_{n+1}=0$, called the hyperplane at infinity. For an affine subvariety $\mathcal V$ of $\Aspace^n\subseteq\PP^n$, the smallest projective subvariety of $\PP^n$ containing $\mathcal V$ is the projective closure of $\mathcal V$, which we denote by $\mathcal V^{\rm cl}$. It is known that $\mathcal V^{\rm cl}\cap\Aspace^n=\mathcal V$, $\dim(\mathcal V^{\rm cl})=\dim(\mathcal V)$, and $\deg(\mathcal V^{\rm cl})=\deg(\mathcal V)$.
\paragraph{Joins of disjoint projective varieties.}
For two distinct points $p,q\in\PP^n$, denote by $\overline{pq}$ the unique projective line passing through them. For two disjoint projective subvarieties $\mathcal V_1,\mathcal V_2\subseteq\PP^n$, define the join $J(\mathcal V_1,\mathcal V_2)$ of $\mathcal V_1$ and $\mathcal V_2$ as
\[J(\mathcal V_1,\mathcal V_2):=\bigcup_{p\in\mathcal V_1,q\in\mathcal V_2}\overline{pq}.\]
\begin{result}{Lemma 2.6 (Harris 1992, Examples 6.17, 11.36, and 18.17)}
$J(\mathcal V_1,\mathcal V_2)$ is a subvariety of $\PP^n$ of dimension $\dim(\mathcal V_1)+\dim(\mathcal V_2)+1$ and degree at most $\deg(\mathcal V_1)\cdot\deg(\mathcal V_2)$.
\end{result}
\begin{result}{Lemma 2.7 (Harris 1992, Exercise 11.6 and Corollary 18.5)}
Let $\mathcal V$ be a nonempty equidimensional subvariety of $\PP^n$ and $H$ a hypersurface of $\PP^n$ not containing an irreducible component of $\mathcal V$. Then $\mathcal V\cap H$ is an equidimensional subvariety of dimension $\dim(\mathcal V)-1$ and degree at most $\deg(\mathcal V)\cdot\deg(H)$ (or an empty set if $\dim(\mathcal V)=0$).
\end{result}
\begin{result}{Corollary 2.8}
Let $\mathcal V$ be a subvariety of $\PP^n$ of dimension $r$, where $0\le r<n$. Then there exists an $(n-r-1)$-subspace $W$ disjoint from $\mathcal V$.
\end{result}
\begin{proof}
It suffices to show that there exist hyperplanes $H_1,\ldots,H_{r+1}$, such that $\mathcal V_i:=\mathcal V\cap(\bigcap_{j=1}^iH_j)$ is empty for some $i\le r+1$. We may inductively choose each $H_i$ such that $C\not\subseteq H_i$ for every irreducible component $C$ of $\mathcal V_{i-1}$, so that $\dim(\mathcal V_i)\le\dim(\mathcal V_{i-1})-1$ by Lemma 2.7. So $\mathcal V_i=\emptyset$ for some $i\le r+1$.
\end{proof}
\begin{result}{Lemma 2.9 (Shafarevich 1994, Section I.6.2, Theorem 6)}
Suppose $\mathcal V_1$ and $\mathcal V_2$ are subvarieties of $\PP^n$ and $\dim(\mathcal V_1)+\dim(\mathcal V_1)\ge n$. Then $\mathcal V_1\cap\mathcal V_2\ne\emptyset$ and $\dim(\mathcal V_1\cap\mathcal V_2)\ge\dim(\mathcal V_1)+\dim(\mathcal V_1)-n$.
\end{result}
\section*{Original Section 3.1: Reducing to the Case of Equidimensional or Irreducible Varieties (journal pp.26--27)}
The following lemma states that to construct $k$-subspace families that are evasive for subvarieties of $\PP^n$, it suffices to consider equidimensional subvarieties of dimension $n-k-1$ (i.e., codimension $k+1$).
\begin{result}{Lemma 3.1}
Let $n,d\in\mathbb N^+$ and $k\in\{0,1,\ldots,n-1\}$. Let $\mathcal F$ be the family of all equidimensional subvarieties of $\PP^n$ of dimension $n-k-1$ and degree at most $d$. Then an $(\mathcal F,\epsilon)$-evasive $k$-subspace family is also $(n,d,\epsilon)$-evasive.
\end{result}
The proof of Lemma 3.1 is based on the following claim.
\begin{result}{Claim 3.2}
Let $\mathcal V$ be an irreducible subvariety of $\PP^n$. There exists a subvariety $\widetilde{\mathcal V}\subseteq\PP^n$ of dimension $n-k-1$ and degree at most $\deg(\mathcal V)$ such that any $k$-subspace of $\PP^n$ that evades $\widetilde{\mathcal V}$ also evades $\mathcal V$.
\end{result}
\begin{proof}
If $\dim(\mathcal V)=n-k-1$, then just let $\widetilde{\mathcal V}=\mathcal V$.
Now assume $\dim(\mathcal V)<n-k-1$. Let $t=(n-k-1)-\dim(\mathcal V)-1$ and let $\widetilde{\mathcal V}$ be the join of $\mathcal V$ and a $t$-subspace disjoint from $\mathcal V$ (which exists by Corollary 2.8). Then $\widetilde{\mathcal V}$ is a projective subvariety of dimension $n-k-1$ and degree at most $\deg(\mathcal V)$ by Lemma 2.6. Suppose $W$ is a $k$-subspace that evades $\widetilde{\mathcal V}$. Then $W$ is disjoint from $\widetilde{\mathcal V}\supseteq\mathcal V$. So $W$ also evades $\mathcal V$.
Finally, assume $\dim(\mathcal V)>n-k-1$. Let $t=\dim(\mathcal V)-(n-k-1)$. By Lemma 2.7, there exist $t$ hyperplanes $H_1,\ldots,H_t$ of $\PP^n$ such that $\mathcal V\cap\bigcap_{i=1}^tH_i$ is equidimensional of dimension $n-k-1$ and degree at most $\deg(\mathcal V)$. Let $\widetilde{\mathcal V}=\mathcal V\cap\bigcap_{i=1}^tH_i$. Suppose $W$ is a $k$-subspace that evades $\widetilde{\mathcal V}$. Then $W\cap\widetilde{\mathcal V}=(W\cap\mathcal V)\cap\bigcap_{i=1}^tH_i=\emptyset$. Again by Lemma 2.7, we have $\dim(W\cap\mathcal V)\le t-1=\dim(\mathcal V)+\dim(W)-n$. So $W$ also evades $\mathcal V$.
\end{proof}
\begin{proof}[Proof of Lemma 3.1]
Consider a projective subvariety $\mathcal V\subseteq\PP^n$ of degree at most $d$. Let $\mathcal V_1,\ldots,\mathcal V_s$ be the irreducible components of $\mathcal V$. For each $i\in[s]$, use Claim 3.2 to choose a projective subvariety $\widetilde{\mathcal V}_i\subseteq\PP^n$ of dimension $n-k-1$ and degree at most $\deg(\mathcal V_i)$ such that any $k$-subspace that evades $\widetilde{\mathcal V}_i$ also evades $\mathcal V_i$. Let $\widetilde{\mathcal V}=\bigcup_{i=1}^s\widetilde{\mathcal V}_i$. Then $\widetilde{\mathcal V}\in\mathcal F$. By construction, any $k$-subspace that evades $\widetilde{\mathcal V}$ also evades $\mathcal V$. It follows that an $(\mathcal F,\epsilon)$-evasive $k$-subspace family is also $(n,d,\epsilon)$-evasive.
\end{proof}
We further reduce to the case of irreducible varieties at the cost of blowing up the parameter $\epsilon$ by a factor of $d$. This is useful as we need irreducibility later in Lemma 3.13.
\begin{result}{Lemma 3.3}
Let $n,d\in\mathbb N^+$ and $k\in\{0,1,\ldots,n-1\}$. Let $\mathcal F'$ be the family of all irreducible subvarieties of $\PP^n$ of dimension $n-k-1$ and degree at most $d$. Then an $(\mathcal F',\epsilon)$-evasive $k$-subspace family is also an $(n,d,d\epsilon)$-evasive $k$-subspace family.
\end{result}
\begin{proof}
Let $\mathcal F$ be as in Lemma 3.1. Each $\mathcal V\in\mathcal F$ has at most $d$ irreducible components, which are all in $\mathcal F'$ since their degrees are bounded by $d$. By definition and the union bound, if a $k$-subspace family $\mathcal H$ is $(\mathcal F',\epsilon)$-evasive, then it is also $(\mathcal F,d\epsilon)$-evasive. Combining this with Lemma 3.1 proves the lemma.
\end{proof}
\section*{Original Chow forms in Stiefel coordinates (journal pp.31--32)}
We may also express the Chow form in Stiefel coordinates, i.e., in the entries of a generating matrix of a linear subspace. This expression has the advantage that it is an actual polynomial rather than a member of the abstract vector space $(R/I)_d$.
Formally, let $A^*$ be a $(k+1)\times(n+1)$ variable matrix whose $(i,j)$-th entry is a variable $Y_{i,j}$. Define the ring homomorphism
\[\phi^\sharp:R=\FF[X_S:S\in\binom{[n+1]}{k+1}]\to\FF[Y_{i,j}:i\in[k+1],j\in[n+1]]\]
that sends each variable $X_S$ to $\det(A^*_{[k+1],S})$. Define the Chow form of $\mathcal V$ in Stiefel coordinates to be
\[\widetilde R_{\mathcal V}:=\phi^\sharp(P_{\mathcal V})\in\FF[Y_{i,j}:i\in[k+1],j\in[n+1]]\]
where $P_{\mathcal V}\in R_d$ is a lift of $R_{\mathcal V}\in(R/I)_d$. Note that $I$ is precisely the kernel of $\phi^\sharp$. So $\widetilde R_{\mathcal V}$ is uniquely determined by $\mathcal V$ up to scalars. By construction, for any $W\in\GG(k+1,n+1)$ and generating matrix $A=(a_{i,j})_{i\in[k+1],j\in[n+1]}$ of $W$, we have $P_{\mathcal V}(\phi(W))=\widetilde R_{\mathcal V}(A):=\widetilde R_{\mathcal V}(a_{1,1},\ldots,a_{k+1,n+1})$. So $\widetilde R_{\mathcal V}$ vanishes at $A$ iff $\PP W\in\GG(k,n)$ intersects $\mathcal V$.
\paragraph{Chow forms in dual Stiefel coordinates.}
Similarly, we may express the Chow form in dual Stiefel coordinates, i.e., in the entries of a parity check matrix of a linear subspace.
More specifically, choose a homogeneous polynomial $Q_{\mathcal V}\in\FF[X_S:S\in\binom{[n+1]}{n-k}]$ that defines the set of $k$-subspaces intersecting $\mathcal V$ in terms of dual Pl\"ucker coordinates. As primal and dual Pl\"ucker coordinates are equivalent, $Q_{\mathcal V}$ can be obtained from the polynomial $P_{\mathcal V}$ above by simply negating and renaming variables. Next, compose $Q_{\mathcal V}$ with a ring homomorphism that substitutes dual Pl\"ucker coordinates with dual Stiefel coordinates. The resulting polynomial, which we denote by $\widetilde R^\vee_{\mathcal V}\in\FF[Y_{i,j}:i\in[n-k],j\in[n+1]]$, is called the Chow form of $\mathcal V$ in dual Stiefel coordinates.
We note that the Chow form $\widetilde R_{\mathcal V}$ in primal Stiefel coordinates is a homogeneous polynomial of degree $(k+1)d$ in $(k+1)(n+1)$ variables, whereas the Chow form $\widetilde R^\vee_{\mathcal V}$ in dual Stiefel coordinates is a homogeneous polynomial of degree $(n-k)d$ in $(n-k)(n+1)$ variables. This suggests that it is more convenient to use the Chow form in primal (resp. dual) Stiefel coordinates when $k$ is small (resp. $n-k$ is small).

\section*{Original Section 3.3: Explicit Constructions of Variety Evasive Subspace Families (journal pp.32--40)}
Let $n,d\in\mathbb N^+$, $k\in\{0,1,\ldots,n-1\}$, and $\epsilon\in(0,1)$. In this subsection, we prove the Main theorem (Theorem 1.7) by constructing explicit projective or affine $k$-subspace families that are $(n,d,\epsilon)$-evasive.
We first prove Theorem 1.7 in the projective case and then derive the affine case from it by viewing $\Aspace^n$ as an open subset of $\PP^n$. For the projective case, we present two constructions. The first one is simple and only uses $\epsilon$-hitting sets for low degree polynomials (Lemma 2.3). But the size of the resulting subspace family is polynomial only when both $d$ and $k$ (or $n-k$) are bounded. Next, we give a more sophisticated construction, which yields subspace families of polynomial size as long as $d$ is bounded.
\subsection*{3.3.1 Simple Construction}
We first present a simple construction of $(n,d,\epsilon)$-evasive $k$-subspace families on $\PP^n$.
First assume $k+1\le n-k$. In this case, construct a $k$-subspace family $\mathcal H$ on $\PP^n$ as follows:
\begin{enumerate}
\item Use Lemma 2.3 to compute an $\epsilon$-hitting set $T$ for the family of polynomials $f\in\FF[Y_{i,j}:i\in[k+1],j\in[n+1]]$ of degree at most $(k+1)d$ such that $|T|=\operatorname{poly}(\binom{(k+1)(n+1+d)}{(k+1)d},1/\epsilon)$. Think of $T$ as a collection of $(k+1)\times(n+1)$ matrices over $\FF$.
\item Initialize $\mathcal H=\emptyset$. For each matrix $A\in T$, if $A$ has full row rank $k+1$, add to $\mathcal H$ the $k$-subspace $W\in\GG(k,n)$ with the generating matrix $A$.
\end{enumerate}
Next, assume $k+1>n-k$. In this case, construct $\mathcal H$ in a similar way, but use parity check matrices instead of generating matrices. Namely, compute an $\epsilon$-hitting set $T$ for the family of polynomials $f\in\FF[Y_{i,j}:i\in[n-k],j\in[n+1]]$ of degree at most $(n-k)d$ such that $|T|=\operatorname{poly}(\binom{(n-k)(n+1+d)}{(n-k)d},1/\epsilon)$. Think of $T$ as a collection of $(n-k)\times(n+1)$ matrices over $\FF$. For each matrix $A\in T$, add to $\mathcal H$ the $k$-subspace $W\in\GG(k,n)$ with the parity check matrix $A$.
This construction does give an $(n,d,\epsilon)$-evasive $k$-subspace family, as stated by the following lemma.
\begin{result}{Lemma 3.11}
The $k$-subspace family $\mathcal H$ constructed above is $(n,d,\epsilon)$-evasive and has size polynomial in
\[\min\left\{\binom{(k+1)(n+1+d)}{(k+1)d},\binom{(n-k)(n+1+d)}{(n-k)d}\right\}\]
and $1/\epsilon$. Moreover, the total time complexity of computing the linear equations defining the $k$-subspaces in $\mathcal H$ is polynomial in $|\mathcal H|$ (and $\log p$, if $\operatorname{char}(\FF)=p>0$).
\end{result}
\begin{proof}
We only show that $\mathcal H$ is $(n,d,\epsilon)$-evasive since the rest of the lemma is obvious from the construction. Let $\mathcal F$ be the family of all equidimensional subvarieties of $\PP^n$ of dimension $n-k-1$ and degree at most $d$. By Lemma 3.1, it suffices to prove that $\mathcal H$ is $(\mathcal F,\epsilon)$-evasive. Consider any $\mathcal V\in\mathcal F$. We want to show that $\mathcal V\cap W=\emptyset$ for all but at most $\epsilon$-fraction of $W\in\mathcal H$.
First assume $k+1\le n-k$, or equivalently, $\binom{(k+1)(n+1+d)}{(k+1)d}\le\binom{(n-k)(n+1+d)}{(n-k)d}$. The Chow form $\widetilde R_{\mathcal V}$ of $\mathcal V$ in Stiefel coordinates is a nonzero homogeneous polynomial in $\FF[Y_{i,j}:i\in[k+1],j\in[n+1]]$ of degree $(k+1)\deg(\mathcal V)\le(k+1)d$. By the choice of $T$, for all but at most $\epsilon$-fraction of $A\in T$, we have $\widetilde R_{\mathcal V}(A)\ne0$, which implies $\mathcal V\cap W=\emptyset$, where $A$ is a generating matrix of $W$.
By construction, $\mathcal H$ is the collection of $k$-subspaces corresponding to the matrices $A\in T$ of full row rank. So we have ignored the matrices that do not have full row rank. But this does not increase the fraction of ``bad'' $W\in\mathcal H$ since if $A$ does not have full row rank, then the maximal minors of $A$ are all zero, and $\widetilde R_{\mathcal V}(A)$ must be zero. It follows that $\mathcal V\cap W=\emptyset$ for all but at most $\epsilon$-fraction of $W\in\mathcal H$, as desired.
Now assume $k+1>n-k$. The proof in this case is similar and we omit the details. The only difference is that we use the Chow form $\widetilde R^\vee_{\mathcal V}$ in dual Stiefel coordinates instead of $\widetilde R_{\mathcal V}$.
\end{proof}
\subsection*{3.3.2 Improved Construction}
Before presenting the improved construction, we first introduce some notions from algebraic geometry.
\paragraph{Projections.}
Suppose $W$ is a $k$-subspace of $\PP^n$, and $\ell_1,\ldots,\ell_{n-k}\in\FF[X_1,\ldots,X_{n+1}]$ are $n-k$ homogeneous linear polynomials such that $W=V(\ell_1,\ldots,\ell_{n-k})$. Then we have a map $\pi:\PP^n\setminus W\to\PP^{n-k-1}$ defined by
\[\pi:x\mapsto(\ell_1(x),\ldots,\ell_{n-k}(x))\]
which is well-defined since $\ell_1,\ldots,\ell_{n-k}$ never simultaneously vanish on $\PP^n\setminus W$. We say $\pi$ is a projection from $\PP^n\setminus W$ to $\PP^{n-k-1}$ and $W$ is its center. Note that if we lift $\pi$ to the linear map $\pi':\Aspace^{n+1}\to\Aspace^{n-k}$ sending $x\in\Aspace^{n+1}$ to $(\ell_1(x),\ldots,\ell_{n-k}(x))\in\Aspace^{n-k}$, then the center $W$ is simply $\PP\ker(\pi')$.
We need the following lemma.
\begin{result}{Lemma 3.12 (Shafarevich 1994, Section I.5.3, Theorem 7)}
Suppose $\pi:\PP^n\setminus W\to\PP^m$ is a projection with center $W$ and $\mathcal V$ is a subvariety of $\PP^n$ disjoint from $W$. Then $\pi$ restricts to a finite morphism from $\mathcal V$ to $\PP^m$.
\end{result}
\paragraph{Nondegenerate varieties.}
For a subvariety $\mathcal V\subseteq\PP^n$, denote by $\operatorname{span}(\mathcal V)$ the smallest projective subspace that contains $\mathcal V$. We say $\mathcal V$ is nondegenerate if it is not contained in a hyperplane of $\PP^n$, or equivalently, $\operatorname{span}(\mathcal V)=\PP^n$.
We need the following fact from algebraic geometry (see, e.g., Eisenbud \& Harris (1987, Proposition 0) or Harris (1992, Corollary 18.12)).
\begin{result}{Lemma 3.13}
The codimension of a nondegenerate irreducible subvariety $\mathcal V$ of $\PP^n$ is at most $\deg(\mathcal V)-1$.
\end{result}
\paragraph{A two-step construction.}
We now give an improved construction of $(n,d,\epsilon)$-evasive $k$-subspace families on $\PP^n$ as follows.
\begin{enumerate}
\item If $k\le d-2$, just use the previous simple construction. So assume $k>d-2$. Let $k'=d-2<k$, $n'=k'+n-k<n$, and $\epsilon_0=\epsilon/(2d)$.
\item Use Corollary 2.5 to construct a collection $\mathcal H_1$ of $(n'+1)\times(n+1)$ matrix over $\FF$ such that $|\mathcal H_1|=\operatorname{poly}(n,d/\epsilon)$ and for every $(n+1)\times(n'+1)$ matrix $M$ over $\FF$ of rank $n'+1$, all but at most $\epsilon_0$-fraction of $B\in U$ satisfies $\operatorname{rank}(BM)=n'+1$.
We abuse the notation and view $\mathcal H_1$ as a collection of linear maps from $\Aspace^{n+1}$ to $\Aspace^{n'+1}$. Then for any linear subspace $W\subseteq\Aspace^{n+1}$ of dimension $n'+1$, we have $\dim(\pi(W))=\dim(W)=n'+1$ for all but at most $\epsilon_0$-fraction of $\pi\in\mathcal H_1$.
\item Construct a collection $\mathcal H_2$ of linear maps from $\Aspace^{n'+1}$ to $\Aspace^{n-k}$ as follows. First assume $d>1$. Use Lemma 3.11 to construct an $(n',d,\epsilon_0)$-evasive $k'$-subspace family $\mathcal H'_2$ on $\PP^{n'}$ of size polynomial in
\[\min\left\{\binom{(k'+1)(n'+1+d)}{(k'+1)d},\binom{(n-k)(n'+1+d)}{(n-k)d}\right\}\]
and $1/\epsilon_0$. For each $k'$-subspace $W\in\mathcal H_2$, compute a surjective linear map $\pi_W:\Aspace^{n'+1}=\Aspace^{n-k+k'+1}\to\Aspace^{n-k}$ such that $W=\PP\ker(\pi_W)$. Let $\mathcal H_2=\{\pi_W:W\in\mathcal H'_2\}$.
If $d=1$, just let $\mathcal H_2$ be the singleton consisting of the identity map on $\Aspace^{n'+1}=\Aspace^{n-k}$.
\item Initialize $\mathcal H=\emptyset$. For each $(\pi_1,\pi_2)\in\mathcal H_1\times\mathcal H_2$, if $\dim(\ker(\pi_2\circ\pi_1))=k+1$, add the $k$-subspace $\PP\ker(\pi_2\circ\pi_1)$ to $\mathcal H$.
\end{enumerate}
We use the construction above to prove the Main theorem (Theorem 1.7) in the projective case. For convenience, we restate it in the following form.
\begin{result}{Theorem 3.14 (Main theorem in the projective case)}
The $k$-subspace family $\mathcal H$ constructed above is $(n,d,\epsilon)$-evasive and has size $\operatorname{poly}(N(k,d,n),n,1/\epsilon)$. Moreover, the total time complexity of computing the linear equations defining the $k$-subspaces in $\mathcal H$ is polynomial in $|\mathcal H|$ (and $\log p$, if $\operatorname{char}(\FF)=p>0$).
\end{result}
\begin{proof}
The theorem follows from Lemma 3.11 if $k\le d-2$. So assume $k>d-2$. We only show that $\mathcal H$ is $(n,d,\epsilon)$-evasive since the rest of the theorem is obvious from the construction.
Let $\mathcal F$ be the family of all irreducible subvarieties of $\PP^n$ of dimension $n-k-1$ and degree at most $d$. By Lemma 3.3, it suffices to prove that $\mathcal H$ is $(\mathcal F,2\epsilon_0)$-evasive. Consider any $\mathcal V\in\mathcal F$. We want to show that $\mathcal V\cap W=\emptyset$ for all but at most $(2\epsilon_0)$-fraction of $W\in\mathcal H$.
By definition, $\mathcal V$ is a nondegenerate irreducible subvariety of $\operatorname{span}(\mathcal V)$. By Lemma 3.13, the codimension of $\mathcal V$ in $\operatorname{span}(\mathcal V)$ is at most $d-1$. Therefore,
\[\dim(\operatorname{span}(\mathcal V))\le\dim(\mathcal V)+d-1=(n-k-1)+(d-1)=n'.\]
Let $\Lambda\subseteq\PP^n$ be an $n'$-subspace that contains $\operatorname{span}(\mathcal V)$. By the choice of $\mathcal H_1$, all but at most $\epsilon_0$-fraction of $\pi_1\in\mathcal H_1$ satisfies $\PP\ker(\pi_1)\cap\Lambda=\emptyset$. Fix $\pi_1$ that satisfies this condition. Then $\PP\ker(\pi_1)$ is disjoint from $\mathcal V\subseteq\Lambda$. By Lemma 3.12, $\pi_1$ induces a finite morphism $\overline\pi_1:\mathcal V\to\PP^{n'}$.
Let $\mathcal V'=\overline\pi_1(\mathcal V)\subseteq\PP^{n'}$. Then $\mathcal V'$ is a projective subvariety of dimension $\dim(\mathcal V)=n-k-1$ and degree at most $d$. (The degree bound follows from, e.g., an inductive application of Mumford (1976, Proposition 5.5).) By the choice of $\mathcal H_2$, all but at most $\epsilon_0$-fraction of $\pi_2\in\mathcal H_2$ satisfies $\PP\ker(\pi_2)\cap\mathcal V'=\emptyset$. Fix $\pi_2$ that satisfies this condition. (If $d=1$ and $\pi_2$ is the identity map, we regard $\PP\ker(\pi_2)$ as an empty set, in which case this condition is also satisfied.) By Lemma 3.12, $\pi_2$ induces a finite morphism $\overline\pi_2:\mathcal V'\to\PP^{n-k-1}$. So we have a finite morphism $\overline\pi_2\circ\overline\pi_1:\mathcal V\to\PP^{n-k-1}$. Note that $\overline\pi_2\circ\overline\pi_1$ is defined by restricting a projection with center $\PP\ker(\pi_2\circ\pi_1)$ to $\mathcal V$. As $\overline\pi_2\circ\overline\pi_1$ is well-defined on $\mathcal V$, its center $\PP\ker(\pi_2\circ\pi_1)$ is disjoint from $\mathcal V$ and this also forces $\dim\ker(\pi_2\circ\pi_1)=k+1$ by Lemma 2.9.
By the above argument and the construction of $\mathcal H$, all but at most $2\epsilon_0$-fraction of the $k$-subspaces in $\mathcal H$ are disjoint from $\mathcal V$, as desired.
\end{proof}
\section*{Original Section 3.3.4: Proof of the Main theorem in the Affine Case (journal pp.39--40)}
We now prove Theorem 1.7 in the affine case. Recall that we may view $\Aspace^n$ as an open subset of $\PP^n$ via the map $(x_1,\ldots,x_n)\mapsto(x_1,\ldots,x_n,1)$. In this way, $\PP^n$ becomes the disjoint union of $\Aspace^n$ and the hyperplane at infinity $H_\infty$ defined by $X_{n+1}=0$.
We use the following lemma to reduce the affine case to the projective case.
\begin{result}{Lemma 3.15}
Let $n,d\in\mathbb N^+$, $k\in\{0,1,\ldots,n-1\}$, and $\epsilon'\in(0,1/2)$. Suppose $\mathcal H$ is an $(n,d,\epsilon')$-evasive $k$-subspace family on $\PP^n$. Then
\[\mathcal H'=\{W\cap\Aspace^n:W\in\mathcal H,W\not\subseteq H_\infty\}\]
is an $(n,d,\epsilon)$-evasive affine $k$-subspace family on $\Aspace^n$, where $\epsilon=\epsilon'/(1-\epsilon')\le2\epsilon'$. Moreover,
\[\mathcal H''=\{W\in\mathcal H:W\not\subseteq H_\infty\}=\{W^{\rm cl}:W\in\mathcal H'\}\]
is an $(n,d,\epsilon)$-evasive $k$-subspace family on $\PP^n$.
\end{result}
\begin{proof}
By $(n,d,\epsilon')$-evasiveness of $\mathcal H$, at most $\epsilon'$-fraction of $W\in\mathcal H$ are fully contained in $H_\infty$. Throwing away those $k$-subspaces fully contained in $H_\infty$ increases the error parameter $\epsilon'$ by at most a factor of $1/(1-\epsilon')$. Therefore, $\mathcal H''=\{W\in\mathcal H:W\not\subseteq H_\infty\}$ is $(n,d,\epsilon)$-evasive. We want to prove that $\mathcal H'=\{W\cap\Aspace^n:W\in\mathcal H''\}$ is also $(n,d,\epsilon)$-evasive.
Consider a subvariety $\mathcal V\subseteq\Aspace^n$ of degree at most $d$. Let $\mathcal V_1,\ldots,\mathcal V_s$ be the irreducible components of $\mathcal V$. The projective closure $\mathcal V^{\rm cl}$ of $\mathcal V$ has the irreducible components $(\mathcal V_1)^{\rm cl},\ldots,(\mathcal V_s)^{\rm cl}$. Consider a $k$-subspace $W\in\mathcal H''$ that evades $\mathcal V^{\rm cl}$. We just need to prove that $W\cap\Aspace^n$ evades $\mathcal V$. This is true since for each $i\in[s]$,
\begin{align*}
\dim((W\cap\Aspace^n)\cap\mathcal V_i)&\le\dim(W\cap(\mathcal V_i)^{\rm cl})\\
&\le\dim(W)+\dim((\mathcal V_i)^{\rm cl})-n\\
&=\dim(W\cap\Aspace^n)+\dim(\mathcal V_i)-n
\end{align*}
where the second inequality holds since $W$ evades $\mathcal V^{\rm cl}$ and the last equality uses the fact $W\not\subseteq H_\infty$.
\end{proof}
The affine case of Theorem 1.7 now follows easily.
\begin{proof}[Proof of Theorem 1.7 in the affine case]
If $k=n$, just choose $\mathcal H=\Aspace^n$. Now assume $k<n$. Construct an $(n,d,\epsilon/2)$-evasive $k$-subspace family $\mathcal H$ on $\PP^n$ using Theorem 3.14. Then
\[\mathcal H':=\{W\cap\Aspace^n:W\in\mathcal H,W\not\subseteq H_\infty\}\]
is an $(n,d,\epsilon)$-evasive affine $k$-subspace family on $\Aspace^n$ by Lemma 3.15. The nonhomogeneous linear equations defining $W\cap\Aspace^n\in\mathcal H'$ can be easily computed from the homogeneous linear equations defining $W\in\mathcal H$ by letting $X_{n+1}=1$.
\end{proof}
The proof of Theorem 1.7 is now complete.

\section*{Original Section 3.3.3: Derandomizing Noether's Normalization Lemma (journal pp.37--39)}
We now prove Theorem 1.12 and Theorem 1.13. For convenience, we restate the theorems below.
\begin{result}{Theorem 1.12}
Let $n,d\in\mathbb N^+$, $r\in\{0,1,\ldots,n\}$, $k=n-r-1$, and $\epsilon\in(0,1)$. There exists an explicit collection $\mathcal L$ of linear maps $\Aspace^{n+1}\to\Aspace^{r+1}$ of size $\operatorname{poly}(N(k,d,n),n,1/\epsilon)$ such that for every subvariety $\mathcal V\subseteq\PP^n$ of dimension $r$ and degree at most $d$, all but at most $\epsilon$-fraction of $\pi\in\mathcal L$ induce a surjective finite morphism from $\mathcal V$ to $\PP^r$.\footnote{Let $N(k,d,n)=1$ when $r=n$ (i.e., $k=-1$). Similarly, in Theorem 1.13, let $N(k,d,n-1)=1$ when $r=n$ or $r=0$ (i.e., $k=-1$ or $k=n-1$).} Moreover, $\mathcal L$ can be computed in time polynomial in $|\mathcal L|$ (and $\log p$, if $\operatorname{char}(\FF)=p>0$).
\end{result}
\begin{proof}
If $r=n$, just use the identity map on $\Aspace^{n+1}=\Aspace^{r+1}$. So assume $r<n$. Use Theorem 3.14 to construct an $(n,d,\epsilon)$-evasive $k$-subspace family $\mathcal H$ on $\PP^n$ of size $\operatorname{poly}(N(k,d,n),n,1/\epsilon)$. For each $k$-subspace $W\in\mathcal H$, compute a surjective linear map $\pi_W:\Aspace^{n+1}\to\Aspace^{r+1}$ such that $W=\PP\ker(\pi_W)$. Let $\mathcal L=\{\pi_W:W\in\mathcal H\}$. Then $\mathcal L$ is a desired collection of linear maps by Lemma 3.12.
\end{proof}
\begin{result}{Theorem 1.13}
Let $n,d\in\mathbb N^+$, $r\in\{0,1,\ldots,n\}$, $k=n-r-1$, and $\epsilon\in(0,1)$. There exists an explicit collection $\mathcal L$ of linear maps $\Aspace^n\to\Aspace^r$ of size $\operatorname{poly}(N(k,d,n-1),n,1/\epsilon)$ such that for every subvariety $\mathcal V\subseteq\Aspace^n$ of dimension $r$ and degree at most $d$, all but at most $\epsilon$-fraction of $\pi\in\mathcal L$ restrict to a surjective finite morphism from $\mathcal V$ to $\Aspace^r$. Moreover, $\mathcal L$ can be computed in time polynomial in $|\mathcal L|$ (and $\log p$, if $\operatorname{char}(\FF)=p>0$).
\end{result}
\begin{proof}
If $r=n$, just use the identity map on $\Aspace^n=\Aspace^r$. If $r=0$, use the only map $\Aspace^n\to\Aspace^0$. So assume $0<r<n$. Regard $\Aspace^n$ as an open subset of $\PP^n$ via $(x_1,\ldots,x_n)\mapsto(x_1,\ldots,x_n,1)$. Similarly, regard $\Aspace^r$ as an open subset of $\PP^r$ via $(x_1,\ldots,x_r)\mapsto(x_1,\ldots,x_r,1)$. Let $H_\infty$ be the hyperplane at infinity of $\PP^n$ defined by $X_{n+1}=0$.
Use Theorem 3.14 to construct an $(n-1,d,\epsilon)$-evasive $k$-subspace family $\mathcal H$ on $H_\infty\cong\PP^{n-1}$ of size $\operatorname{poly}(N(k,d,n-1),n,1/\epsilon)$. For each $W\in\mathcal H$, choose $n-k=r+1$ homogeneous linear polynomials $\ell_1,\ldots,\ell_{r+1}\in\FF[X_1,\ldots,X_{n+1}]$ such that $\ell_{r+1}=X_{n+1}$, $\ell_1,\ldots,\ell_r\in\FF[X_1,\ldots,X_n]$, and $W=V(\ell_1,\ldots,\ell_{r+1})$. This is possible as $W\subseteq H_\infty=V(X_{n+1})$.
These $r+1$ linear polynomials determine a projection $\pi_W:\PP^n\setminus W\to\PP^r$, defined by
\begin{align*}
x=(x_1,\ldots,x_{n+1})&\mapsto(\ell_1(x),\ldots,\ell_{r+1}(x))\\
&=(\ell_1(x),\ldots,\ell_r(x),x_{n+1}).
\end{align*}
As $x_{n+1}=1$ for $x\in\Aspace^n$, we have $\pi_W(\Aspace^n)\subseteq\Aspace^r$. Restricting $\pi_W$ on $\Aspace^n$ yields a map $\pi_W|_{\Aspace^n}:\Aspace^n\to\Aspace^r$, which is a linear map as $\ell_1,\ldots,\ell_r$ are homogeneous linear polynomials in $\FF[X_1,\ldots,X_n]$. Let $\mathcal L=\{\pi_W|_{\Aspace^n}:W\in\mathcal H\}$.
Let $\mathcal V$ be a subvariety of $\Aspace^n$ of dimension $r$ and degree at most $d$. Its projective closure $\mathcal V^{\rm cl}$ has dimension $\dim(\mathcal V)=r$ and degree $\deg(\mathcal V)\le d$. By the definition of $\mathcal V^{\rm cl}$, none of the irreducible components of $\mathcal V^{\rm cl}$ is fully contained in $H_\infty$. So by Lemma 2.7, the projective subvariety $\mathcal V^{\rm cl}\cap H_\infty$ has dimension $r-1$ and degree at most $d$.
By the choice of $\mathcal H$, all but at most $\epsilon$-fraction of $W\in\mathcal H$ are disjoint from $\mathcal V^{\rm cl}\cap H_\infty$ and hence from $\mathcal V^{\rm cl}$. So we just need to prove that for every $W\in\mathcal H$ disjoint from $\mathcal V^{\rm cl}$ and the corresponding projection $\pi_W$, the map $\pi_W|_{\mathcal V}:\mathcal V\to\Aspace^r$ is a surjective finite morphism. This follows from Lemma 3.12 and the fact that $\mathcal V=\mathcal V^{\rm cl}\cap(\pi_W)^{-1}(\Aspace^r)$.
\end{proof}

\begin{thebibliography}{99}
\bibitem{journalref1} Manindra Agrawal, Rohit Gurjar, Arpita Korwar \& Nitin Saxena (2015). Hitting-sets for ROABP and sum of set-multilinear circuits. SIAM Journal on Computing 44(3), 669–697.
\bibitem{journalref2} Manindra Agrawal, Chandan Saha, Ramprasad Saptharishi \& Nitin Saxena (2016). Jacobian hits circuits: Hitting sets, lower bounds for depth-D occur-k formulas and depth-3 transcendence degreek circuits. SIAM Journal on Computing 45(4), 1533–1562.
\bibitem{journalref3} Manindra Agrawal \&VVinay (2008). Arithmetic circuits: A chasm at depth four. In Proceedings of the 49th Annual IEEE Symposium on Foundations of Computer Science, 67–75. IEEE.
\bibitem{journalref4} Pablo Azcue (1992). On the dimension of the Chow varieties. Ph.D. thesis, Harvard University.
\bibitem{journalref5} Malte Beecken, Johannes Mittmann \& Nitin Saxena (2013). Algebraic independence and blackbox identity testing. Information and Computation 222, 2–19.
\bibitem{journalref6} Markus Bläser \& Anurag Pandey (2020). Polynomial Identity Testing for Low Degree Polynomials with Optimal Randomness. In Approximation, Randomization, and Combinatorial Optimization. Algorithms and Techniques (APPROX/RANDOM), 8:1–8:13. Schloss Dagstuhl – Leibniz-Zentrum für Informatik.
\bibitem{journalref7} Andrej Bogdanov (2005). Pseudorandom generators for low degree polynomials. In Proceedings of the 37th Annual ACM Symposium on Theory of Computing, 21–30. ACM.
\bibitem{journalref8} Nader Bshouty (2014). Testers and their applications. In Proceedings of the 5th Conference on Innovations in Theoretical Computer Science, 327–352. ACM.
\bibitem{journalref9} Fabrizio Catanese (1992). Chow varieties, Hilbert schemes, and moduli spaces of surfaces of general type. Journal of Algebraic Geometry 1(4), 561–595.
\bibitem{journalref10} Arthur Cayley (1860). On a new analytical representation of curves in space. The Quarterly Journal of Pure and Applied Mathematics 3, 225–236.
\bibitem{journalref11} Arkadev Chattopadhyay, Ankit Garg \& Suhail Sherif (2021). Toward Stronger Counterexamples to the Log-Approximate-Rank Conjecture. In 41st IARCS Annual Conference on Foundations of Software Technology and Theoretical Computer Science (FSTTCS 2021), 13:1– 13:16. Schloss Dagstuhl – Leibniz-Zentrum für Informatik.
\bibitem{journalref12} Wei-Liang Chow \& B.L. van der Waerden (1937). Zur algebraischen Geometrie. IX. Mathematische Annalen 113(1), 692–704.
\bibitem{journalref13} Gil Cohen \& Amnon Ta-Shma (2013). Pseudorandom Generators for Low Degree Polynomials from Algebraic Geometry Codes. In Electronic Colloquium on Computational Complexity (ECCC), volume 20, 155.
\bibitem{journalref14} John Dalbec \& Bernd Sturmfels (1995). Introduction to Chow forms. Invariant Methods in Discrete and Computational Geometry 37–58.
\bibitem{journalref15} Richard A. Demillo \& Richard J. Lipton (1978). A probabilistic remark on algebraic program testing. Information Processing Letters 7(4), 193–195.
\bibitem{journalref16} Thomas W Dubé (1993). A combinatorial proof of the effective Nullstellensatz. Journal of Symbolic Computation 15(3), 277–296.
\bibitem{journalref17} Pranjal Dutta, Prateek Dwivedi \& Nitin Saxena (2021). Deterministic Identity Testing Paradigms for Bounded Top-Fanin Depth-4 Circuits. In 36th Computational Complexity Conference (CCC 2021), 11:1–11:27. Schloss Dagstuhl – Leibniz-Zentrum für Informatik.
\bibitem{journalref18} Zeev Dvir (2012). Extractors for varieties. Computational Complexity 21(4), 515–572.
\bibitem{journalref19} Zeev Dvir, Ariel Gabizon \& Avi Wigderson (2009). Extractors and rank extractors for polynomial sources. Computational Complexity 18(1), 1–58.
\bibitem{journalref20} Zeev Dvir, János Kollár \& Shachar Lovett (2014). Variety evasive sets. Computational Complexity 23(4), 509–529.
\bibitem{journalref21} Zeev Dvir \& Shachar Lovett (2012). Subspace evasive sets. In Proceedings of the 44th Annual ACM Symposium on Theory of Computing, 351–358. ACM.
\bibitem{journalref22} Zeev Dvir \& Amir Shpilka (2007). Locally decodable codes with two queries and polynomial identity testing for depth 3 circuits. SIAM Journal on Computing 36(5), 1404–1434.
\bibitem{journalref23} David Eisenbud (1995). Commutative Algebra: with a View Toward Algebraic Geometry. Springer-Verlag.
\bibitem{journalref24} David Eisenbud \& Joe Harris (1987). On varieties of minimal degree. In Proceedings of Symposia in Pure Mathematics, volume 46, 3–13. American Mathematical Society.
\bibitem{journalref25} David Eisenbud \& Joe Harris (1992). The dimension of the Chow variety of curves. Compositio Mathematica 83(3), 291–310.
\bibitem{journalref26} Michael A. Forbes (2014). Polynomial identity testing of read-once oblivious algebraic branching programs. Ph.D. thesis, Massachusetts Institute of Technology.
\bibitem{journalref27} Michael A. Forbes \& Venkatesan Guruswami (2015). Dimension Expanders via Rank Condensers. In Approximation, Randomization, and Combinatorial Optimization. Algorithms and Techniques (APPROX/RANDOM). Schloss Dagstuhl – Leibniz-Zentrum für Informatik.
\bibitem{journalref28} Michael A. Forbes, Ramprasad Saptharishi \& Amir Shpilka (2014). Hitting sets for multilinear read-once algebraic branching programs, in any order. In Proceedings of the 46th Annual ACM Symposium on Theory of Computing, 867–875. ACM.
\bibitem{journalref29} Michael A. Forbes \& Amir Shpilka (2012). On identity testing of tensors, low-rank recovery and compressed sensing. In Proceedings of the 44th Annual ACM Symposium on Theory of Computing, 163–172. ACM.
\bibitem{journalref30} Michael A. Forbes \& Amir Shpilka (2013). Explicit Noether normalization for simultaneous conjugation via polynomial identity testing. In Approximation, Randomization, and Combinatorial Optimization. Algorithms and Techniques (APPROX/RANDOM), 527–542. Springer Berlin Heidelberg.
\bibitem{journalref31} Michael A. Forbes \& Amir Shpilka (2018). A PSPACE construction of a hitting set for the closure of small algebraic circuits. In Proceedings of the 50th Annual ACM Symposium on Theory of Computing, 1180–1192. ACM.
\bibitem{journalref32} William Fulton (1997). Young Tableaux: With Applications to Representation Theory and Geometry, volume 35. Cambridge University Press.
\bibitem{journalref33} Ariel Gabizon \& Ran Raz (2008). Deterministic extractors for affine sources over large fields. Combinatorica 28(4), 415–440.
\bibitem{journalref34} Israel M. Gelfand, Mikhail M. Kapranov \& Andrei V. Zelevinsky (1994). Discriminants, Resultants and Multidimensional Determinants. Birkhäuser.
\bibitem{journalref35} Marc Giusti, Klemens Hägele, Grégoire Lecerf, Joël Marchand \& Bruno Salvy (2000). The projective Noether Maple package: computing the dimension of a projective variety. Journal of Symbolic Computation 30(3), 291–307.
\bibitem{journalref36} Marc Giusti \& Joos Heintz (1993). La détermination des points isolés et de la dimension d’une variété algébrique peut se faire en temps polynomial. Computational Algebraic Geometry and Commutative Algebra (Cortona, 1991) 216–256.
\bibitem{journalref37} Mark L. Green \& Ian Morrison (1986). The equations defining Chow varieties. Duke Mathematical Journal 53(3), 733–747.
\bibitem{journalref38} Lucio Guerra (1999). Complexity of Chow varieties and number of morphisms on surfaces of general type. Manuscripta Mathematica 98(1), 1–8.
\bibitem{journalref39} Zeyu Guo (2021). Variety Evasive Subspace Families. In 36th Computational Complexity Conference (CCC 2021), 20:1–20:33. Schloss Dagstuhl – Leibniz-Zentrum für Informatik.
\bibitem{journalref40} Zeyu Guo \& Noga Ron-Zewi (2021). Efficient List-Decoding With Constant Alphabet and List Sizes. IEEE Transactions on Information Theory 68(3), 1663–1682.
\bibitem{journalref41} Zeyu Guo, Nitin Saxena \& Amit Sinhababu (2019). Algebraic Dependencies and PSPACE Algorithms in Approximative Complexity over Any Field. Theory of Computing 15(16), 1–30.
\bibitem{journalref42} Ankit Gupta (2014). Algebraic Geometric Techniques for Depth-4 PIT \& Sylvester-Gallai Conjectures for Varieties. In Electronic Colloquium on Computational Complexity (ECCC), volume 21, 130.
\bibitem{journalref43} Venkatesan Guruswami \& Swastik Kopparty (2016). Explicit subspace designs. Combinatorica 36(2), 161–185.
\bibitem{journalref44} Venkatesan Guruswami, Nicolas Resch \& Chaoping Xing (2021). Lossless dimension expanders via linearized polynomials and subspace designs. Combinatorica 1–35.
\bibitem{journalref45} Venkatesan Guruswami, Carol Wang \& Chaoping Xing (2016). Explicit list-decodable rank-metric and subspace codes via subspace designs. IEEE Transactions on Information Theory 62(5), 2707–2718.
\bibitem{journalref46} Venkatesan Guruswami \& Chaoping Xing (2013). List decoding Reed-Solomon, Algebraic-Geometric, and Gabidulin subcodes up to the Singleton bound. In Proceedings of the 45th Annual ACM Symposium on Theory of Computing, 843–852. ACM.
\bibitem{journalref47} Venkatesan Guruswami, Chaoping Xing \& Chen Yuan (2018). Subspace designs based on algebraic function fields. Transactions of the American Mathematical Society 370(12), 8757–8775.
\bibitem{journalref48} Joe Harris (1992). Algebraic Geometry: A First Course. Springer- Verlag.
\bibitem{journalref49} Joos Heintz (1983). Definability and fast quantifier elimination in algebraically closed fields. Theoretical Computer Science 24(3), 239– 277.
\bibitem{journalref50} Joos Heintz \& Claus-Peter Schnorr (1980). Testing polynomials which are easy to compute (Extended Abstract). In Proceedings of the 12th annual ACM Symposium on Theory of Computing, 262–272. ACM.
\bibitem{journalref51} Joos Heintz \& Malte Sieveking (1981). Absolute primality of polynomials is decidable in random polynomial time in the number of variables. In International Colloquium on Automata, Languages, and Programming, 16–28. Springer Berlin Heidelberg.
\bibitem{journalref52} David Hilbert (1890). Ueber die Theorie der algebraischen Formen. Mathematische Annalen 36(4), 473–534.
\bibitem{journalref53} Gabriela Jeronimo, Teresa Krick, Juan Sabia \& Martín Sombra (2004). The computational complexity of the Chow form. Foundations of Computational Mathematics 4(1), 41–117.
\bibitem{journalref54} Michael Joswig \& Thorsten Theobald (2013). Polyhedral and Algebraic Methods in Computational Geometry. Springer Science \& Business Media.
\bibitem{journalref55} Zohar S. Karnin \& Amir Shpilka (2011). Black box polynomial identity testing of generalized depth-3 arithmetic circuits with bounded top fan-in. Combinatorica 31(3), 333.
\bibitem{journalref56} Steven Kleiman (1976). Problem 15. Rigorous foundation of Schubert’s enumerative calculus. In Proceedings of Symposia in Pure Mathematics, volume 28, 445–482.
\bibitem{journalref57} Adam R. Klivans \& Daniel Spielman (2001). Randomness efficient identity testing of multivariate polynomials. In Proceedings of the 33rd Annual ACM Symposium on Theory of Computing, 216–223. ACM.
\bibitem{journalref58} János Kollár (1988). Sharp effective Nullstellensatz. Journal of the American Mathematical Society 963–975.
\bibitem{journalref59} János Kollár (1996). Rational Curves on Algebraic Varieties. Springer.
\bibitem{journalref60} Swastik Kopparty, Noga Ron-Zewi, Shubhangi Saraf \& Mary Wootters (2023). Improved list decoding of folded Reed-Solomon and multiplicity codes. SIAM Journal on Computing 52(3), 794–840.
\bibitem{journalref61} Teresa Krick (2002). Straight-line programs in polynomial equation solving. Foundations of Computational Mathematics 312, 96–136.
\bibitem{journalref62} Joseph M. Landsberg (2012). Tensors: Geometry and Applications. Graduate studies in mathematics. American Mathematical Society.
\bibitem{journalref63} Joseph M. Landsberg (2015). Geometric complexity theory: an introduction for geometers. Annali dell’universita’di Ferrara 61(1), 65– 117.
\bibitem{journalref64} Brian Lehmann (2017). Asymptotic behavior of the dimension of the Chow variety. Advances in Mathematics 308, 815–835.
\bibitem{journalref65} Hendrik W. Lenstra, Jr. (1990). Algorithms for finite fields. Number theory and cryptography 154, 76–85.
\bibitem{journalref66} Nutan Limaye, Srikanth Srinivasan \& Sébastien Tavenas (2024). Superpolynomial lower bounds against low-depth algebraic circuits. Communications of the ACM 67(2), 101–108.
\bibitem{journalref67} Chi-Jen Lu (2012). Hitting set generators for sparse polynomials over any finite fields. In Proceedings of the 27th Conference on Computational Complexity, 280–286. IEEE.
\bibitem{journalref68} Partha Mukhopadhyay (2016). Depth-4 identity testing and Noether’s normalization lemma. In Proceedings of the 11th International Computer Science Symposium in Russia, 309–323. Springer International Publishing.
\bibitem{journalref69} Ketan Mulmuley (2017). Geometric complexity theory V: Efficient algorithms for Noether normalization. Journal of the American Mathematical Society 30(1), 225–309.
\bibitem{journalref70} David Mumford (1976). Algebraic Geometry I: Complex Projective Varieties. Springer-Verlag.
\bibitem{journalref71} Masayoshi Nagata (1962). Local rings. Interscience Tracts in Pure and Applied Mathematics .
\bibitem{journalref72} Emmy Noether (1926). Der Endlichkeitssatz der Invarianten endlicher linearer Gruppen der Charakteristik p. Nachrichten von der Gesellschaft der Wissenschaften zu Göttingen, Mathematisch- Physikalische Klasse 1926, 28–35.
\bibitem{journalref73} Luis M. Pardo \& Daniel Sebastián (2022). A promenade through correct test sequences I: Degree of constructible sets, Bézout’s Inequality and density. Journal of Complexity 68, 101 588.
\bibitem{journalref74} Shir Peleg \& Amir Shpilka (2021). Polynomial time deterministic identity testing algorithm for $\Sigma$[3]$\Pi$$\Sigma$$\Pi$[2] circuits via Edelstein–Kelly type theorem for quadratic polynomials. In Proceedings of the 53rd Annual ACM Symposium on Theory of Computing, 259–271. ACM.
\bibitem{journalref75} Shir Peleg \& Amir Shpilka (2022). A generalized Sylvester–Gallaitype theorem for quadratic polynomials. Forum of Mathematics, Sigma 10, e112.
\bibitem{journalref76} Nitin Saxena (2009). Progress on Polynomial Identity Testing. Bulletin of the EATCS 99, 49–79.
\bibitem{journalref77} Nitin Saxena (2013). Progress on Polynomial Identity Testing - II. Electronic Colloquium on Computational Complexity (ECCC) 20, 186. URL http://eccc.hpi-web.de/report/2013/186 .
\bibitem{journalref78} Nitin Saxena \& Comandur Seshadhri (2012). Blackbox identity testing for bounded top-fanin depth-3 circuits: The field doesn’t matter. SIAM Journal on Computing 41(5), 1285–1298.
\bibitem{journalref79} Jacob T. Schwartz (1980). Fast probabilistic algorithms for verification of polynomial identities. Journal of the ACM 27(4), 701–717.
\bibitem{journalref80} Igor R. Shafarevich (1994). Basic Algebraic Geometry 1: Varieties in Projective Space. Springer-Verlag.
\bibitem{journalref81} Amir Shpilka (2020). Sylvester-Gallai type theorems for quadratic polynomials. Discrete Analysis .
\bibitem{journalref82} Amir Shpilka \& Amir Yehudayoff (2010). Arithmetic Circuits: A Survey of Recent Results and Open Questions. Now Publishers Inc.
\bibitem{journalref83} Richard Zippel (1979). Probabilistic algorithms for sparse polynomials. In International Symposium on Symbolic and Algebraic Manipulation, 216–226. Springer Berlin Heidelberg.
\end{thebibliography}
\end{document}
